% DRAFT (LLM-drafted from the claims ledger; the author reviews, edits and confirms every claim).
\section{Introduction}
\label{sec:intro}

Many applied classification problems reduce to a decision among a known set of options: route a customer message to
one of 77 banking intents, pick the tool an agent should call, or choose the next step in a workflow. A recent class
of systems targets exactly this setting. We call them \emph{decision models}: given an instruction, a closed list of
options (optionally with descriptions) and optional evidence such as labelled examples, they return a probability
distribution over the options. Some are offered as closed APIs that advertise calibrated probabilities
\citep{typesafe-2026-jev}; others are released with open weights \authorcheck{cite the Clef and pplx-decider model
cards}. All of them report results on public benchmarks.

Those results are hard to compare. A decision model's accuracy depends not only on the model but on everything that
goes into the request: the instruction, the wording of the options and, in the few-shot setting, which labelled
examples are retrieved and in what order. Published numbers are usually produced with each system's own request
format and retrieval, on different subsets, and with metrics whose definitions differ (binning of calibration error,
rounding of probabilities). A difference between two published numbers can therefore reflect the evidence the
systems were given rather than the systems.

This paper proposes and applies an \textbf{equal-evidence protocol} for comparing decision models
(\S\ref{sec:protocol}). Every system receives byte-identical requests whose hashes are recorded before scoring; every
scored run is pre-registered in a public commit before it starts; systems are compared item by item with exact
McNemar tests and Holm correction within declared families; probabilities are compared with paired bootstrap
intervals; and a model scored on any engine other than its release code must first pass a per-item parity gate
against that code. The protocol is deliberately simple. Its value lies in what it exposes when it is applied
strictly.

We apply it to one closed API (\jev), three open decision models (\clef, \flash{} and \pplx) and two low-cost open
baselines that we built as a testbed (\S\ref{sec:systems}): a training-free zero-shot jury of three sub-1B
encoders (\gj) and a retrieval-augmented 24-shot stack around a 4B reader (\orb). Our main case is Banking77
\citep{casanueva-etal-2020-efficient} at 24 shots, where a public reproduction of the closed API provides the
reference request format, retrieved examples and per-item predictions \citep{smith-2026-jev-banking77}.
Two further pre-registered blocks test breadth: a second 24-shot dataset (Bitext customer support) on identical
evidence for all five systems (\S\ref{sec:res-bitext}), and thirteen code and workflow decision tasks for the three
open decision models (\S\ref{sec:res-13task}).

\paragraph{Findings.}
\begin{enumerate}
  \item \textbf{A closed API's published result can be replicated exactly.} Re-sending the reproduction's
  \fsN{} request bodies to the pinned API version gives \fsJevAAcc\%, equal to the published \fsJevPubAcc\%
  (b/c $=\fsJevAVsJevB/\fsJevAVsJevC$, \fsJevAAgreement\% identical choices; \S\ref{sec:res-replication}).
  \item \textbf{The retrieved evidence matters as much as the model.} With our retrieval rule instead of the
  reproduction's, the same API gains \fsJevBMinusJevAPt{} points ($p=\fsJevBVsJevAP$). Our open baseline beat the
  published API result ($p=\fsCleanVsJevP$), but at equal retrieval the two tie ($p=\fsCleanVsJevBP$), and under the
  reproduction's rule our baseline also ties it ($p=\fsParBestVsJevP$). A post hoc factorial analysis attributes the
  gain mainly to the absence of a per-class cap in retrieval. The three open decision models gain little or nothing
  from the same change; none of their changes is significant (\S\ref{sec:res-retrieval}).
  \item \textbf{Accuracy and calibration rank models differently.} On identical evidence the most accurate model
  (\flash, \fsFlashAcc\%) has the highest calibration error (ECE \fsFlashEce), while a 27B open model (\pplx)
  matches the best ECE and has the lowest Brier score. Both calibration patterns hold under either retrieval rule,
  although on our retrieval the three open models' accuracies cannot be separated (\S\ref{sec:res-calibration}).
  \item \textbf{Metric and reproducibility pitfalls change conclusions.} Rounding of returned probabilities to
  two decimals makes log loss unusable against one API (\S\ref{sec:find-rounding}); platform-dependent tie-breaking
  in BM25 changes the retrieved examples for \fsStableSetDiff{} of \fsN{} test items yet only \fsStableChanged{}
  predictions (\S\ref{sec:find-ties}); and a configuration chosen on the test set reached a
  significance that the pre-registered one does not need (\S\ref{sec:find-selection}).
  \item \textbf{Zero-shot comparisons need training-data disclosure.} Decision models lead zero-shot by wide
  margins, but the open models either list benchmark datasets in their training data or do not disclose it, and
  one of our own baseline members inherits benchmark training data through its initialisation
  (\S\ref{sec:find-exposure}).
\end{enumerate}

\paragraph{Contributions.} (i) The equal-evidence protocol, with the tooling to apply it (request manifests,
parity gate, pre-registration chain); (ii) the measured findings above; (iii) a release of every request body,
per-item prediction and analysis script, so that new decision models can be added to the comparison on the same
evidence (\S\ref{sec:release}); and (iv) two documented, low-cost open baselines that serve as reproducible
reference points rather than as competitors: on equal evidence, the open decision models are more accurate
(\S\ref{sec:res-calibration}).
